\section{Software design and assurance}

The stable runtime is organized around validated domain objects, canonical
schemas, numerical kernels, diagnostics, provenance, and adapters. The Rust
workspace contains the authoritative numerical implementation and a narrow,
versioned C application binary interface with explicit ownership and error
transport. The Python package
uses PyO3 and Maturin for its native extension and retains the facade and
workflow responsibilities needed by users. R and Julia consume the same
interface;
their packages do not reimplement VOI policy. Optional dependencies such as
JAX, plotting libraries, and web integrations are isolated from the base core.

Correctness is supported by language-neutral fixtures, unit and integration
tests, property-based and fuzz tests, metamorphic checks, mutation testing,
coverage gates, Rust diagnostics, ABI checks, benchmark regression budgets,
and clean-install tests. The release process produces platform wheels, a
source archive, checksums, a software bill of materials, provenance evidence,
and a signed GitHub release. Documentation is maintained in Astro/Starlight.
The v1.0.0 release is available from GitHub and the Python Package Index, while
external registry and
curation outcomes remain tracked separately.

The manuscript follows software-citation guidance~\cite{smith2016software}.
Its canonical source is authored LaTeX. The repository builds a deterministic
source-only archive, audits the review PDF, compiles against the supported
arXiv TeX Live generations, and independently converts the source to a
semantic web document. These checks establish repository readiness, not submission or
acceptance by arXiv or a journal.
